\chapter*{GENERAL CONCLUSION}
\markboth{\MakeUppercase{GENERAL CONCLUSION}}{}
\addcontentsline{toc}{chapter}{GENERAL CONCLUSION}
\adjustmtc
\thispagestyle{MyStyle}

[Placeholder] This concluding chapter synthesizes the contributions of the work and reflects on what was achieved, why it matters, and what comes next.

[Placeholder — paragraph one: what was achieved] This paragraph will summarize the concrete deliverables of the project: a working RFQ lifecycle management service, a trust-bound conversational copilot grounded in platform evidence, an integrated platform that exposes both engines through a coherent user experience, and a validation strategy that combines backend rigor with formal copilot metrics and golden scenario testing.

[Placeholder — paragraph two: why it matters] This paragraph will articulate the business and engineering significance of the work, situating Itqan against the audit pain points identified at the client site and against the broader challenge of building trustworthy AI systems for industrial estimation contexts.

[Placeholder — paragraph three: technical skills and engineering judgment demonstrated] This paragraph will identify the technical competencies exercised during the project: microservices design, the FastAPI and SQLAlchemy stack, LLM integration with epistemic guardrails, the BACAB layered pattern, structured project management with Scrumban and ADR practices, and the discipline of strategic reframing when the original problem definition proved infeasible.

[Placeholder — paragraph four: innovation contribution] This paragraph will name the central innovation contribution of the work: the trust-boundary architecture as a defensible pattern for constrained agentic systems, separating language understanding (owned by the LLM) from operational truth, access control, evidence selection, and response boundaries (owned by deterministic platform components).

[Placeholder — paragraph five: personal perspective and future directions] This paragraph will close on a personal reflection: what the project taught about engineering judgment in the face of ambiguity, what remains to be built across the four-pillar vision, and how the platform might evolve toward multi-client deployment and federated learning across RFQ histories.